% =====================================================================
% §4 Diagnosis-Preserving Enhancement (C1 — hardest main leg). SHORT (9-page) version.
% Serves Claim C1; ACCEPTANCE C1.1-C1.5. caveman OFF.
% Anti-drift locks per STORY: deterministic restoration not generative; FiLM positioned for
%   diagnostic fidelity not pixel quality; completeness partially recoverable (R1);
%   Prop3 carried by E7 plus non-vacuity; derive/sketch wording (R2). Full proofs -> Appendix A2.
% Macros assumed in preamble: \visienhance \Tw \qbar \Ldp \KL \Zenh \Zref \tauenh \tauhigh
% =====================================================================

\section{Diagnosis-Preserving Enhancement}\label{sec:enhance}

Of the five quality dimensions, four (sharpness, brightness, colour temperature, and contrast) admit
deterministic restoration, whereas completeness ($q_3$, a truncated field of view) is only
\emph{partially recoverable}: enhancement cannot synthesise tissue that was never captured, so the agent
treats low completeness as a query-for-retake signal rather than an enhancement target. This split ---
which dimensions are recoverable, and within what severity window --- is the boundary this section
formalises and \cref{sec:exp} measures.

\subsection{Architecture: NAFNet with quality-conditional FiLM}\label{sec:arch}
\visienhance{} is a NAFNet backbone \citep{chen2022simple} with FiLM modulation \citep{perez2018film}
conditioned on the per-input quality vector, so restoration strength adapts to \emph{how} an image is
degraded rather than applying a fixed correction. The restoration map $\Tw(x,q)$ is
\emph{deterministic}: we avoid generative restoration (diffusion or adversarial backbones), because such
models can hallucinate diagnostic structure --- inventing a pigment network or vascular pattern that was
never present --- which is unacceptable clinically (\cref{sec:discussion} analyses why diffusion fails
here). FiLM is \emph{not} a fidelity mechanism: its empirical role (E8, \cref{sec:exp}) is
quality-conditional \emph{diagnostic} fidelity, not perceptual sharpening, and removing it leaves
per-image PSNR essentially unchanged.

\subsection{DP-Loss: a feature-level diagnosis-preserving objective}\label{sec:dploss}
Training minimises
\begin{equation}
\mathcal{L}=\mathcal{L}_1
  +\lambda_{\mathrm{LPIPS}}\,\mathcal{L}_{\mathrm{LPIPS}}
  +\lambda_{\mathrm{DP}}\,\Ldp
  +\lambda_{\mathrm{q}}\,\mathcal{L}_{\mathrm{hinge}},
\label{eq:total-loss}
\end{equation}
where $\mathcal{L}_1$ and $\mathcal{L}_{\mathrm{LPIPS}}$ govern pixel and perceptual fidelity, and the
diagnosis-preserving term is a \emph{feature-level} KL,
\begin{equation}
\Ldp \;=\; \KL\!\big(p_\phi^{\mathrm{enh}} \,\big\|\, p_\phi^{\mathrm{ref}}\big),
\label{eq:dploss}
\end{equation}
computed in the latent space of a frozen diagnostic feature extractor (the $1536\times7\times7$
penultimate features of an independent EfficientNet-B3 probe). The point is that $\Ldp$ aligns the
enhanced image with its clean reference \emph{in the space that drives the diagnosis}, not in pixel
space: a restorer can be pixel-faithful yet still shift the diagnostic posterior, and \eqref{eq:dploss}
penalises exactly that movement. The quality hinge $\mathcal{L}_{\mathrm{hinge}}$ encourages
$\qbar(\Tw(x,q))>\qbar(x)$ within the salvage band, tying restoration to a measured quality gain.

\subsection{Mutual-information preservation}\label{sec:prop3lemma3}
% Lemma stated in body; directional entropy Proposition (prop:entropy) deferred to Appendix A2. R2: derive/sketch, not "prove".
We connect the DP objective to diagnosis through a feature-level mutual-information lower bound, stated
here and derived in full in Appendix~\ref{app:a2}.
% Compact statements of Prop 3 + Lemma 3 for the main body (§4.4).
% Full proofs: Appendix A2.1 / A2.2 (appendix/A2_prop3_lemma3.tex).

\begin{proposition}[Enhancement Entropy Reduction]\label{prop:entropy}
Under Assumptions A1--A4 (quality improvement, calibrated scoring, monotone prior,
diagnostic preservation), for $x$ in the salvage band and $q'=f_\eta(\Tw(x,q))$,
\begin{equation}
\E_{x}\!\big[\Hent(\hat{p}_{\Tw(x,q)})\big] \;\leq\;
\E_{x}\!\big[\Hent(\hat{p}_x)\big]
- \tfrac{d}{2}\log\tfrac{\sigma^2(\qbar)}{\sigma^2(\qbar')} + 2\beta\sqrt{\epsilon},
\label{eq:prop3}
\end{equation}
where the first negative term is the \qvib{} entropy gap (Lemma~1, strictly positive
since $\qbar'>\qbar \Rightarrow \sigma^2$ shrinks) and the residual $2\beta\sqrt{\epsilon}$
is the diagnostic-preservation slack (Lemma~\ref{lem:dp}).
\end{proposition}

\begin{lemma}[Diagnosis-Preserving Mutual-Information Bound]\label{lem:dp}
Let $q_\theta$ be $\Lq$-Lipschitz in total variation and let the label entropy be
bounded by $M=\log K$. If $\Ldp := \KL(p_\phi^{\mathrm{enh}} \,\|\, p_\phi^{\mathrm{ref}}) \leq \epsilon$, then
\begin{equation}
I(\Zenh; Y) \;\geq\; I(\Zref; Y) - \beta\sqrt{\epsilon},
\qquad \beta = \frac{M\,\Lq}{\sqrt{2}}.
\label{eq:lemma3}
\end{equation}
The $\sqrt{\epsilon}$ rate is Pinsker-optimal; a linear-$\epsilon$ bound would require a
stronger $\chi^2$-divergence assumption (Appendix~A2.2.5). The proof combines Pinsker's
inequality with a Lipschitz coupling argument (Appendix~A2.2).
\end{lemma}

In words, Lemma~\ref{lem:dp} says that driving the feature-level KL small provably caps how much
diagnostic information enhancement can destroy; its empirical counterpart is experiment~E7
(\cref{sec:exp}), which is the evidence we rely on. A companion directional entropy result
(Proposition~\ref{prop:entropy}, Appendix~\ref{app:a2}) characterises only the sign of one variance
term and is \emph{not} a net entropy reduction: a quality-conditioned probe finds that predictive
entropy does not fall after enhancement. The information-preservation role is therefore carried by
Lemma~\ref{lem:dp} (validated by E7) together with the Stage-1 PSNR$\,\geq30$\,dB non-vacuity
condition. We state these as derivations under the listed assumptions rather than as watertight proofs.

\subsection{Three-stage training protocol}\label{sec:training}
We train $\visienhance$ in three stages, each isolating one term in \eqref{eq:total-loss}: \textbf{(1)}
pixel restoration ($\mathcal{L}_1$ + LPIPS) trained to the PSNR gate that also guards the non-vacuity
condition of \eqref{eq:lemma3}; \textbf{(2)} DP-Loss fine-tuning that tightens $\Ldp$ toward the budget
$\epsilon$ --- the stage whose ablation is E7; and \textbf{(3)} a quality-hinge stage calibrating the
$\qbar$ gain inside the salvage band $[\tauenh,\tauhigh]$, so enhancement is credited only where it
measurably raises quality.
